\documentclass[a4paper,10pt,twoside]{report}

\usepackage{Packages}

\addto\captionsbritish{\renewcommand{\contentsname}{Tabla de Contenidos}}

%\linenumbers

\newcommand{\shortdoctitle}{Doc title}
\newcommand{\doctitle}{Formalización del problema de la partición del grafo}

\newcommand{\authorone}{Laura Rodríguez Navas \\ rodrigueznavas@posgrado.uimp.es}
\newcommand{\monthYear}{Junio 2020}

\author{\me}

\hypersetup
{
    pdfauthor   = {\authorone},
    pdftitle    = {\shortdoctitle},
    pdfsubject  = {\doctitle}
}

\begin{document}

\begin{titlepage}
\begin{center}
\includegraphics[height=2.1cm]{Figures/uimp-logo.png} \\
\LARGE
Resolución de problemas con metaheurísticos \\
\Large
Informe \\
\large

\vspace*{10cm}

\setlength{\TPHorizModule}{1mm}
\setlength{\TPVertModule}{\TPHorizModule}

\newlength{\backupparindent}
\setlength{\backupparindent}{\parindent}
\setlength{\parindent}{0mm}

\begin{textblock}{155}(32,89)
    \vspace*{1mm}
    \huge
    \textbf{\doctitle \\}
    \Large
    \vspace*{5mm}
    % \textit{\docsubtitle} \\
    \vspace*{5mm}
    \Large
     \begin{tabular}{c c c}
            \authorone
    \end{tabular} \\
\end{textblock}

\vfill
%\docdate \\
\large
\monthYear \\

\setlength{\parindent}{\backupparindent}
\end{center}
\end{titlepage} 

\normalsize

\chapter*{Resumen}\label{chapter:Resumen}
\setcounter{page}{0}
\pagenumbering{arabic}
\addcontentsline{toc}{chapter}{Resumen}
Dentro del amplio campo de la Teoría de Grafos, en este trabajo se considera el problema de la partición del grafo: dado un grafo no dirigido cuyo número de vértices es par y cuyas aristas tienen asociado un peso, se trata de encontrar la partición del conjunto de vértices en dos subconjuntos de tamaño similar de manera que se minimice la suma de los pesos de aquellas aristas que unen vértices que están en diferentes conjuntos

En el capítulo inicial se presenta una breve introducción al problema. En el segundo capítulo, se proponen tres codificaciones para este problema, señalando sus aspectos más relevantes (positivos y negativos), además de ejemplos. Concretamente, las codificaciones son implementaciones en Python 3 de los algoritmos basados en técnicas metaheurísticas: Kernighan-Lin, Spectral Bisection y Multilevel Spectral Bisection. Y en el tercer capítulo, se concluye con una comparativa entre las diferentes codificaciones.

Las codificaciones se encuentran en el repositorio:

\begin{center}
	\url{https://github.com/lrodrin/masterAI/tree/master/A2}\label{GitHub}
\end{center}


\tableofcontents

\chapter{Introducción a la Partición de Grafos}\label{chapter:Introducción}
El problema de dividir un grafo en diversas particiones del mismo tamaño tiene muchas utilidades en el área de la electrónica, así como en el desarrollo de sistemas operativos. Este problema ha sido demostrado como un problema NP-completo (\url{https://es.wikipedia.org/wiki/NP-hard}), lo que implica que las soluciones para él no pueden ser encontradas en tiempos razonables. El particionamiento de grafos es un problema importante que tiene aplicaciones extensas en muchas áreas.

\chapter{Algoritmos basados en metaheurísticas}\label{chapter:Algoritmos}
%Dado que el problema de la partición de un grafo es un problema NP-completo, las soluciones prácticas se basan en la heurística y existen dos categorías: local y global. Los métodos locales más conocidos son el algoritmo Kernighan-Lin\cite{KernighanLin} y el algoritmo Fiduccia-Mattheyses\cite{FiducciaMattheyses}. Su principal inconveniente es la partición inicial arbitraria del conjunto de vértices, que puede afectar la calidad de la solución óptima final. En cambio, los métodos globales se basan en las propiedades de todo el grafo y no se basan en una partición inicial arbitraria. El ejemplo más común es la partición espectral, donde una partición se deriva de los vectores propios aproximados de la matriz de adyacencia, o la agrupación espectral que agrupa los vértices del grafo usando la descomposición propia de la matriz Laplaciana del grafo.

\section{Kernighan-Lin}
%El algoritmo de Kernighan-Lin\cite{KernighanLin} es un algoritmo heurístico para resolver el problema de la partición de un grafo con complejidad computacional $O({n}^2logn)$. Este algoritmo, propuesto en 1970 por Shen Lin y Brian Kernighan, tiene aplicaciones importantes para el diseño de circuitos digitales y VLSI. 
%
%Es un ejemplo de Algoritmo Voraz.
%
%Comienza con alguna partición que satisfaga el requisito de tamaño que intercambie repetidamente nodos entre las particiones.

%1
Kernighan-Lin es un método para dividir un grafo que contiene nodos y vértices en subconjuntos separados que están conectados de manera óptima. Dado que se puede usar un grafo para representar una red eléctrica que contiene bloques, el algoritmo Kernighan-Lin se puede extender a la división de circuitos en subcircuitos.

Una aplicación importante del algoritmo es en los circuitos VLSI\cite{KernighanLin}\cite{Ravikumar}. Concretamente se usa para encontrar un número mínimo de conexiones entre particiones para mejorar la velocidad o disminuir el consumo de energía.

Kernighan-Lin es un algoritmo:

\begin{itemize}
	\setlength{\parskip}{0pt}
	\setlength{\itemsep}{0pt plus 1pt}
	\item \textbf{Iterativo}. Esto significa que el grafo/circuito inicialmente ya está particionado, pero la aplicación de Kernighan-Lin intentará mejorar u optimizar la partición inicial. Kernighan-Lin es iterativo en lugar de constructivo.
	\item \textbf{Voraz}. Esto significa que el algoritmo hará cambios si hay un beneficio inmediato sin considerar otras formas posibles de obtener una solución óptima.
	\item \textbf{Determinista} porque se obtendrá el mismo resultado cada vez que se aplique el algoritmo. El mismo resultado será el mismo número de redes que cruzan la bisección, pero no necesariamente las mismas redes.
\end{itemize}


%2
El algoritmo Kernigan-Lin\cite{KernighanLin}, a menudo abreviado como K-L, es uno de los primeros algoritmos de partición de grafos y fue desarrollado originalmente para optimizar la colocación de circuitos electrónicos en tarjetas de circuito impreso para minimizar el número de conexiones entre las tarjetas.

El algoritmo K-L no crea particiones, sino que las mejora iterativamente. La idea original era tomar una partición aleatoria y aplicarle Kernigan-Lin. Esto se repetiría varias veces y se elegiría el mejor resultado. Mientras que para grafos pequeños esto ofrece resultados razonables, es bastante ineficiente para tamaños de problemas más grandes.

Hoy en día, el algoritmo se usa para mejorar las particiones encontradas por otros algoritmos. Como veremos, K-L tiene una visión algo \textit{"local"} del problema, tratando de mejorar las particiones mediante el intercambio de nodos vecinos. Por lo tanto, complementa muy bien los algoritmos que tienen una visión más \textit{"global"} del problema, pero tienden a ignorar las características locales. Un ejemplo de este tipo de algoritmos seria la partición espectral (ver sección \ref{spectral_bisection}).

C. Fiduccia y R. Mattheyses realizaron importantes avances prácticos en \cite{FiducciaMattheyses} quienes mejoraron el algoritmo de tal manera que una iteración se ejecuta en $O(n)$ en lugar de $O({n}^2 \, log \, n)$. A continuación, primero se describe el algoritmo original y después se discuten estas y otras mejoras.

\subsection{Descripción}

\subsection{Mejoras}

\subsection{Ejemplo}


\section{Spectral Bisection}\label{spectral_bisection}

\subsection{Descripción}

\subsection{Ejemplo}

\section{Multilevel Spectral Bisection}

\subsection{Descripción}

\subsection{Ejemplo}

\section{Comparativa entre los diferentes algoritmos}



\chapter{Comparativa entre los diferentes algoritmos}\label{chapter:Comparativa}
En esta sección, veremos una comparativa entre los algoritmos que hemos descrito en las secciones previas, tras ser aplicados sobre diferentes grafos aleatorios, también conocidos como grafos de Erdös y Renyi\cite{ErdosRenyi} o grafos binomiales. Pero antes de ver la comparativa, en primer lugar, definimos el concepto de grafo aleatorio.

Se dice que un grafo es aleatorio si la presencia y colocación de sus aristas siguen una distribución aleatoria. Por tanto, un grafo aleatorio no puede diseñarse bajo ningún criterio concreto. Los grafos aleatorios también son conocidos como grafos de Erdös y Renyi porqué utilizan el modelo de Erdös y Renyi. El nombre de este modelo proviene de los matemáticos Paul Erdős y Alfréd Rényi, quienes lo presentaron por primera en 1959.

El modelo de Erdös y Renyi (a veces abreviado como modelo ER), es el método que se ha empleado para la generación de los grafos aleatorios de prueba que se usan en la comparativa de las codificaciones. Para ello, los grafos se han construido mediante la conexión de los nodos al azar. Cada arista se incluye en el grafo con una probabilidad de \textit{p} independiente de las otras aristas del grafo. De manera equivalente, todos los grafos con \textit{n} nodos y \textit{m} aristas tienen la misma probabilidad. Esa probabilidad es:

\begin{center}
	$p^m(1 - p)^{{n \choose 2} - m}$
\end{center}

En la implementación de los grafos aleatorios se ha utilizado la librería de Python:  \href{https://networkx.github.io/documentation/latest/reference/generated/networkx.generators.random\_graphs.erdos\_renyi\_graph.html?highlight=erdos\%20renyi#networkx.generators.random\_graphs.erdos\_renyi\_graph}{NetworkX}. Concretamente el método que se ha utilizado nos devuelve un grafo no dirigido, donde se elige cada una de las aristas posibles con probabilidad \textit{p} = 0.7, según el número de nodos de entrada \textit{n}. En particular, el caso \textit{p} = 0.7 se corresponde con el caso en el los \textit{n} vértices se eligen con mayor probabilidad (\textit{p} < 0.5). Los pesos de las aristas también se han generado aleatoriamente dentro de un intervalo de pesos de 1 a 20.

Podemos observar la codificación acabada de describir, que se ejecuta en tiempo $O$($n^2$).

\lstset{language=Python}    
\begin{lstlisting}[frame=single]  
G = nx.erdos_renyi_graph(n, 0.7)

for u, v in G.edges():
  if u != v:
    G[u][v]['label'] = random.randrange(1, 20)
\end{lstlisting}

Y en la tabla siguiente, se presenta una comparación entre los algoritmos codificados sobre los grafos aleatorios generados en este informe, en términos de tiempo computacional. Se muestra el tiempo (en segundos) en completar las ejecuciones de los algoritmos sobre los grafos aleatorios para distintos números de vértices \textit{n}. Para ello se ha utilizado un MacBook Pro, con un procesador de cuatro núcleos de 2.8 GHz y con 16 GB de memoria RAM. Es de esperar que, a mayor número de vértices, los algoritmos tengan mayor tiempo de ejecución.

\renewcommand{\tablename}{Tabla}
\begin{table}[h]	
	\begin{center}
		\begin{tabular}{|c|c|c|c|}
			\hline
			\textit{n} & 307 & 552 & 861 \\
			\hline
			Kernighan-Lin & 0.0117 & 0.0214 & 0.0451\\
			\hline
			Spectral Bisection & 0.0021 & 0.0031 & 0.0060 \\
			\hline
			MSB & 0.0025 & 0.0031 & 0.0037 \\ 
			\hline
		\end{tabular}
		\vspace{1mm}
		\caption{Tabla comparativa de los algoritmos.}
	\end{center}
\end{table}

Podemos ver como los algoritmos Kernighan-Lin, Spectral Bisection y Multilevel Spectral Bisection (MSB) tienen unos tiempos de ejecución similares cuando el número de vértices \textit{n} es más pequeño. Fenómeno que cambia cuando casi se duplican el número de vértices a 552. 

\newpage
Además, si nos fijamos en el algoritmo Kernighan-Lin, cuando el número de vértices casi se duplican a 552, el tiempo de ejecución del algoritmo también se duplica. Parece que el algoritmo va aumentando el tiempo de ejecución a medida que aumenta el número de vértices de manera proporcionada.

En cambio, el algoritmo Spectral Bisection obtiene unos resultados muy parecidos con el algoritmo Multilevel Spectral Bisection (MSB), pero cuando el número de vértices empieza a ser muy elevado, la eficiencia de este algoritmo baja drásticamente. Que no es el caso del algoritmo Multilevel Spectral Bisection (MSB).

Como conclusión general de esta breve comparación entre los algoritmos podemos decir que el algoritmo Multilevel Spectral Bisection (MSB) es el más eficiente, con tiempos de ejecución más pequeños, para los grafos aleatorios no dirigidos que se ha codificado, y el algoritmo Kernighan-Lin es el algoritmo menos eficiente, con tiempos de ejecución más grandes. También podemos observar que los tiempos de ejecución no son muy grandes, ya que solo se ha medido el tiempo de ejecución de los algoritmos y además el número de vértices totales no es muy grande. Aun así, nos ha permitido ver las diferencias en tiempo de ejecución de las tres codificaciones que se han implementado en este informe.

La elección del tiempo de ejecución frente a la calidad de las particiones para la comparación de los algoritmos ha dependido del hecho que el problema de partición de grafos sea un problema NP-completo\cite{NPCompleteness}. Que recordemos que esto quiere decir que los algoritmos puede que no nos ofrezcan la solución óptima en cada ejecución, pero sí que nos dan una buena solución al menos la mayor parte del tiempo. O sea, que como es difícil comparar la calidad de las particiones entre los diferentes algoritmos se elijó la medida del tiempo computacional utilizado para cada algoritmo. Por otro lado, en el contexto de matriz-vector de los algoritmos que utilizan la bisección espectral, también solo nos interesaba el tiempo total. 

En el caso de los algoritmos Spectral Bisection y Multilevel Spectral Bisection (MSB), podría ser que, si solo utilizamos una determinada matriz una vez, el algoritmo rápido que entrega solo particiones de calidad media en general podría ser más rápido que el algoritmo más lento con particiones de mejor calidad. Pero si usamos la misma matriz (o diferentes matrices con el mismo grafo) a menudo, el algoritmo más lento podría ser preferible. De hecho, hay incluso más factores a considerar. Hasta ahora queríamos que las diferentes particiones tuvieran el mismo tamaño. Como hemos visto en la sección 2.2, podría ser una ventaja aceptar particiones de un tamaño ligeramente diferente para lograr una mejor solución. Todo esto debería demostrar que no existe un mejor algoritmo único para todas las situaciones y que los diferentes algoritmos descritos en las anteriores secciones pueden tener diferentes aplicaciones.

Aunque en general, los algoritmos de particionamiento complejos tardan más en calcular las particiones. Sin embargo, crean particiones que a menudo son más equilibradas y/o minimizan las conexiones entre particiones son mejores que los algoritmos de partición menos complejos. Ya que, el objetivo es reducir el tiempo de cálculo de una ejecución. En consecuencia, en algunos casos un algoritmo menos complejo sería preferible a uno más complejo. También nos encontramos con que la eficiencia cambia con el número de particiones y que los algoritmos producen resultados diferentes dependiendo de los atributos del grafo de entrada.

\chapter{Conclusiones}\label{chapter:Conclusiones}
Hemos utilizado el mismo grafo aleatorio para los tres ejemplos descritos de las codificaciones de los algoritmos.

El algoritmo Multilevel Spectral Bisection es el más eficiente que se ha codificado, y el algoritmo Kernighan-Lin es el menos eficiente.

Para la comparativa solo se ha tenido en cuenta el tiempo de ejecución de los algoritmos. No se ha añadido el tiempo de ejecución empleado en la creación de los grafos, antes y después de la partición.

Cada vez que ejecutamos un algoritmo las soluciones cambian. Así que, por eso, los hemos comparado por tiempo de ejecución. Ya que, al ser algoritmos metaheurísticos, todas las soluciones pueden ser óptimas y válidas. Eso quiere decir, que, aunque un algoritmo sea más eficiente, no obtendrá la solución óptima.

Como hemos visto, a menudo existe una relación entre el tiempo de ejecución y la calidad de la solución. Algunos algoritmos funcionan muy rápido pero solo encuentran una solución de medio calidad mientras que otros tardan mucho tiempo pero ofrecen soluciones excelentes e incluso otros se puede sintonizar entre ambos extremos.

La elección del tiempo frente a la calidad.

Las particiones pequeñas en general podrían ser más rápidas que un algoritmo más lento con mejor calidad
particiones Pero si usamos la misma matriz (o diferentes matrices con el mismo gráfico)
a menudo, el algoritmo más lento podría ser preferible. tamaño de corte
Todo esto debería demostrar que no existe un mejor algoritmo único para todas las situaciones.
y que los diferentes algoritmos descritos en los siguientes capítulos tienen todos sus
aplicaciones.

\renewcommand{\bibname}{Bibliografia}
\bibliographystyle{unsrt}
\bibliography{Bibliografia}

\end{document}